% Modernised reading — fonds Alexandre Grothendieck, shelfmark 151, batch 2
% (pages 21-40).
%
% This is an INTERPRETATION, not a transcription. It re-reads the pages in
% today's notation and vocabulary, states the hypotheses the manuscript leaves
% implicit, and drops the critical apparatus so that the mathematics reads
% straight through. Everything the brackets and underlines carried — a doubtful
% reading, a notation he used differently, a missing passage — moves into a
% footnote.
%
% It is held to being mathematically correct as it stands, not merely faithful
% to the page. Where the manuscript is loose or mistaken, the true statement is
% what appears here, and the footnote says what the page has. In any dispute of
% substance the transcription governs: whoever cites Grothendieck must cite
% batch-02.fr.tex.
%
% Scope: the folder was read WHOLE before any of this was written — all four
% batches, pages 1-74, in one pass — and one file is written per batch, this
% being the second. The batch stops mid-sentence at the foot of page 40; the
% sentence is completed on page 41, in batch 3, and what it was going to say is
% known here because the whole folder was read first.
%
% Pass: Opus 5 (claude-opus-5), 2026-08-16 — first pass, unchecked against the
% pages by a human.
\documentclass[11pt,a4paper]{article}
% Le préambule commun à toutes les transcriptions du fonds Grothendieck.
%
% Il tient en une idée : **l'incertitude se compose, elle ne se résout pas.**
% Une transcription qui lisse un mot douteux a détruit la seule chose pour
% laquelle on la faisait — un lecteur doit pouvoir savoir, à chaque endroit,
% ce qui était sur le feuillet et ce qui a été ajouté. Les six macros qui
% suivent sont donc l'appareil critique, et non de la décoration.
%
% Les mêmes commandes sont reconnues par `scripts/render.mjs`, qui produit la
% vue de lecture du site. Une macro ajoutée ici doit l'être aussi là-bas,
% faute de quoi le rendu échouera bruyamment — ce qui est voulu : un rendu
% silencieusement faux est pire qu'un rendu absent.

\ProvidesPackage{grothendieck}[2026/08/08 Transcriptions du fonds Grothendieck]

\usepackage{amsmath,amssymb,amsthm}
\usepackage{mathrsfs}
\usepackage{stmaryrd}
% Les diagrammes commutatifs sont partout dans ce fonds : c'est la langue
% même de la géométrie algébrique de Grothendieck, et une transcription qui
% les rendrait en prose ne transcrirait rien.
\usepackage{tikz-cd}
% Français par défaut — c'est la langue des feuillets — mais l'anglais est
% chargé aussi : la traduction et le résumé sont des documents anglais, et la
% typographie française (espaces avant les deux-points) y serait une faute.
% Ces documents basculent par \selectlanguage{english} en tête de corps.
\usepackage[english,french]{babel}
\usepackage{fontspec}
\usepackage{geometry}
\geometry{a4paper,margin=25mm,marginparwidth=28mm}
\usepackage{marginnote}
\usepackage{xcolor}
% ulem, not soul. `soul`'s \st and \ul are built on a character-by-character
% scan that XeLaTeX's font machinery breaks: the document fails with an
% undefined control sequence rather than degrading. ulem does the same two jobs
% and is Unicode-engine safe. `normalem` stops it hijacking \emph, which the
% transcriptions use for Grothendieck's own emphasis.
\usepackage[normalem]{ulem}
\usepackage{hyperref}
\hypersetup{colorlinks=true,linkcolor=black,urlcolor=black,pdfborder={0 0 0}}

% --- Métadonnées du lot -----------------------------------------------------
% Reprises telles quelles par le script de rendu, qui les affiche en tête de
% la vue de lecture. Elles fixent ce que le fichier prétend couvrir : sans
% elles, une transcription flotte sans point d'attache dans un fonds de seize
% mille feuillets.

\newcommand{\folder}[1]{\gdef\grothFolder{#1}}
\newcommand{\batch}[1]{\gdef\grothBatch{#1}}
\newcommand{\leaves}[2]{\gdef\grothFirst{#1}\gdef\grothLast{#2}}
\newcommand{\foldertitle}[1]{\gdef\grothFoldertitle{#1}}
\gdef\grothFolder{?}\gdef\grothBatch{?}\gdef\grothFirst{?}\gdef\grothLast{?}\gdef\grothFoldertitle{}

% --- Le résumé --------------------------------------------------------------
%
% Il ouvre la lecture modernisée, sous le titre « Résumé », et il est composé
% en petit. Ce n'est pas un ornement : c'est le seul passage du document qui
% ne soit pas des mathématiques, et un lecteur doit pouvoir le reconnaître
% avant de l'avoir lu — puis le sauter s'il connaît déjà le sujet, ou s'y
% arrêter s'il ne le connaît pas. Le filet qui le suit marque la reprise du
% corps.

\newenvironment{resume}
  {\small\setlength{\parskip}{0.45em}}
  {\par\medskip\hrule\medskip}

% --- Mots-clés ---------------------------------------------------------------
%
% En anglais, en clôture du résumé de la lecture modernisée : le vocabulaire
% moderne sous lequel chercher ce que les feuillets construisent. Le manifeste
% du site (`npm run manifest`) les extrait pour étiqueter la cote entière —
% cette ligne est la source unique des tags, il n'y a pas de fichier à côté.
\newcommand{\keywords}[1]{\par\smallskip\noindent
  {\footnotesize\textsc{Keywords} --- \textit{#1}}\par}

% --- L'appareil critique ----------------------------------------------------

% Le numéro de feuillet, en marge. C'est l'ancre qui permet de régler un
% désaccord sur une formule en regardant *un* feuillet plutôt qu'une chemise
% de six cents. Le site s'en sert aussi pour tourner les pages du fac-similé
% quand on fait défiler la transcription.
\newcommand{\leaf}[1]{%
  \marginnote{\footnotesize\color{gray!70}\textsf{#1}}%
  \label{leaf:#1}\ignorespaces}

% Illisible : on ne devine pas. Un mot inventé qui se lit comme les autres est
% le pire résultat possible.
\newcommand{\ill}{\textcolor{red!60!black}{[\,\ldots\,]}}

% Lecture douteuse : le mot est donné, et signalé comme incertain.
\newcommand{\uncertain}[1]{\uline{#1}}

% Ajout de l'éditeur — une lettre manquante, un mot sous-entendu.
\newcommand{\add}[1]{\textcolor{blue!50!black}{[#1]}}

% Biffé par Grothendieck lui-même. Ce qu'il a rayé fait partie du document :
% une rature dit souvent où la pensée a changé de direction.
\newcommand{\struck}[1]{\sout{#1}}

% Note du transcripteur, sur l'état matériel du feuillet.
\newcommand{\note}[1]{\par\smallskip{\small\color{gray!60!black}\textit{[#1]}}\par\smallskip}

% Note marginale de Grothendieck, distincte du corps du texte.
\newcommand{\marginal}[1]{\par\smallskip\noindent\hspace{1em}%
  {\small\color{gray!60!black}#1}\par\smallskip}

% --- Titre ------------------------------------------------------------------

\AtBeginDocument{%
  \begin{center}
    {\large\bfseries Fonds Alexandre Grothendieck --- cote n\textsuperscript{o}~\grothFolder}\\[2pt]
    {\itshape\grothFoldertitle}\\[4pt]
    {\small feuillets \grothFirst--\grothLast\ (lot \grothBatch)}\\[2pt]
    {\footnotesize\color{gray!60!black}Transcription établie d'après le fac-similé numérisé
     par l'Université de Montpellier.\\
     Les lectures incertaines sont soulignées, les illisibles notées [\,\ldots\,],
     les ajouts entre crochets.}
  \end{center}
  \vspace{1em}}

\endinput


\folder{151}
\batch{2}
\pages{21}{40}
\dating{[à partir de 1981-à partir de 1990]}
\watermark{Édition de démonstration\\interprétation personnelle de l'œuvre}
\foldertitle{Topos stratifié (1981) : notes manuscrites (s.d.) — lecture modernisée}

\begin{document}

\section*{Résumé}

\begin{resume}
Le dossier poursuit une question : un espace découpé en morceaux — une surface
et les courbes qu'elle porte, une variété et son lieu singulier — est-il
déterminé par ses morceaux et par la façon dont chacun s'approche des autres ?
Le texte « Structures stratifiées », dont le sommaire ouvrait le lot précédent,
commence ici, à la page 26, et il commence par le cas le plus simple : deux
morceaux seulement. Un fermé $Y$ dans un espace $X$, et ce qui reste,
$X^{*} = X \setminus Y$.

L'obstacle est immédiat et vaut d'être vu sur un dessin. Un cône, sa pointe, et
le reste : la pointe est fermée, le reste est ouvert, leur intersection est
vide, et rien dans ces deux données ne dit que le cône se referme. Il faut donc
enregistrer, en plus des deux morceaux, la manière dont l'un s'approche de
l'autre. C'est ce que fait le \emph{tube} $\mathcal{V}_{Y,X}$ — un voisinage
infinitésimal de $Y$ dans $X$, qui se rétracte sur $Y$ — et son \emph{tube
épointé} $\mathcal{V}^{*}_{Y,X}$, ce qu'il en reste quand on ôte $Y$. Ces deux
espaces, avec $Y$ et $X^{*}$, sont les seuls dont le texte se servira, ici et
partout dans la suite du dossier.

L'espace est alors recollé : $X$ est la somme amalgamée de $\mathcal{V}_{Y,X}$
et de $X^{*}$ le long de $\mathcal{V}^{*}_{Y,X}$. De là, un théorème de van
Kampen — le groupe fondamental de $X$ est le produit amalgamé de ceux de $Y$ et
de $X^{*}$ au-dessus de celui du tube épointé —, une suite exacte d'homotopie où
apparaissent deux groupes que Grothendieck baptise, en empruntant à
l'arithmétique, \emph{inertie} et \emph{décomposition}, et une description des
faisceaux qui sont localement constants sur chaque morceau : ce sont des
triplets, un système local de chaque côté et un morphisme de comparaison entre
eux au-dessus du tube épointé.

Puis le piège, et il occupe les trois dernières pages. Tout ce qui précède ne
retient des morceaux que leur groupe fondamental. Est-ce assez pour retrouver le
type d'homotopie de $X$ ? Grothendieck prend l'exemple le plus simple qu'on
puisse imaginer — la sphère $\mathbb{P}^{1}_{\mathbb{C}}$, c'est-à-dire le plan
complexe plus un point à l'infini, dont les morceaux sont un point, un plan et
un cercle, tous les trois aussi simples que possible — et trouve que la donnée
obtenue est \emph{homotopiquement triviale}, alors que la sphère ne l'est pas.
Le niveau du groupe fondamental ne suffit donc pas, et le texte s'arrête au
milieu de la phrase qui va dire pourquoi : il faut monter d'un cran, jusqu'aux
gerbes. La suite est au lot suivant.

Il vaut la peine de noter comment il traverse ce moment : « J'avais craint un
moment, par cet exemple, que le type d'homotopie de $X$ ne pouvait pas se
récupérer […]. Mais je me suis convaincu ensuite que cette crainte n'était pas
fondée. » Ce qui échoue n'est pas le programme, c'est un invariant trop grossier,
et la distinction est faite en une ligne.

Les cinq premières pages du lot n'appartiennent pas à ce texte : elles achèvent
les feuilles de calcul du lot précédent, sur un système à trois étages muni des
permutations des coordonnées, et elles y apportent une réponse complète dans le
cas pointé.

\keywords{van Kampen theorem, tubular neighbourhood, homotopy pushout, constructible sheaf, exit-path category}
\end{resume}

\pagerange{21}{24}
\section*{Le système à trois étages : le cas pointé, achevé}

Les quatre feuilles numérotées 4 à 7 par l'auteur poursuivent, sans coupure, le
calcul des pages 2 à 8 et 16 à 20. Le système est celui de deux ensembles reliés
par un couple section-rétraction,
\[
  d \colon E_{1} \longrightarrow E_{2}, \qquad
  p \colon E_{2} \longrightarrow E_{1}, \qquad p\,d = \mathrm{id}_{E_{1}},
\]
avec $\mathfrak{S}_{2} = \{1,\sigma\}$ opérant sur $E_{1}$, $\mathfrak{S}_{3}$
sur $E_{2}$, un point marqué $e_{1} \in E_{1}$ fixé par $\sigma$, et
$e_{2} = d(e_{1})$.\footnote{L'auteur écrit ici $X_{2}$ et $X_{3}$ pour $E_{1}$
et $E_{2}$, et $X_{1} = \{e_{1}\}$ pour le point ; ces $X$ n'ont aucun rapport
avec les $X_{i}$ de « Structures stratifiées », qui commencent cinq pages plus
loin. Nous gardons la numérotation par degré, $E_{n}$, dans tout le dossier.}
C'est la page 21 qui donne enfin les conventions en coordonnées, valables pour
tout le dossier :
\[
  d(x,y) = (x,y,y), \qquad p(x,y,z) = (x,y), \qquad
  \rho(x,y,z) = (y,z,x),
\]
\[
  \sigma_{1}(x,y,z) = (x,z,y), \qquad \sigma_{2}(x,y,z) = (z,y,x), \qquad
  \sigma_{3}(x,y,z) = (y,x,z),
\]
chaque $\sigma_{i}$ étant donc la transposition qui \emph{fixe} la $i$-ième
lettre.

\subsection*{Une contradiction apparente, et sa levée}

La page encadre trois relations,
\[
  p\,d = \mathrm{id}_{E_{1}}, \qquad p\,\rho\,d = e, \qquad
  p\,\rho^{2}d = \sigma,
\]
où $e$ désigne l'application constante de valeur $e_{1}$ ; et deux relations
d'équivariance, $\sigma_{1}d = d$ et $p\,\sigma_{3} = \sigma\,p$. La seconde
paraît démentie par les coordonnées de la même page, qui donnent
$p\,\rho\,d(x,y) = (y,y)$, lequel n'est pas constant.

Il n'y a pas de contradiction, et la levée est instructive. Dans le modèle en
coordonnées, $E_{n} = F^{n+1}$ pour un ensemble $F$, on a $(y,y) =
\delta_{0}p_{1}(x,y)$ : c'est la diagonale, non une constante. La relation
générale, valable sans hypothèse, est
\[
  p\,\rho\,d = \delta_{0}\,p_{1},
\]
et elle ne devient $p\,\rho\,d = e$ que sous l'hypothèse posée en tête de page,
$E_{0}$ réduit à un point — auquel cas $\delta_{0}p_{1}$ se factorise par un
point et est bien constante. Les deux écritures de la page appartiennent à deux
régimes : les coordonnées au modèle général, les relations encadrées au cas
pointé.\footnote{La transcription signale le désaccord et ne le tranche pas,
comme il se doit. C'est ici, où l'on peut écrire les deux régimes côte à côte,
qu'il se tranche. La même remarque vaut pour la troisième page du lot
précédent, où $p_{12}\rho\delta_{1} = \delta_{0}p_{0}$ est écrit dans le cas
général et « application constante » dans le cas pointé.}

\subsection*{Le sous-système engendré par $d(E_{1})$}

Soit $E'_{2}$ le sous-$\mathfrak{S}_{3}$-ensemble de $E_{2}$ engendré par
$d(E_{1})$. Le stabilisateur de $d(\xi)$ est $\{1,\sigma_{1}\}$ pour tout
$\xi \neq e_{1}$, et $\mathfrak{S}_{3}$ tout entier pour $\xi = e_{1}$ ; d'où
une orbite ponctuelle et des orbites à trois éléments, en bijection avec
$E^{*}_{1} = E_{1} \setminus \{e_{1}\}$ :
\[
  E'_{2} \;\simeq\; \{e_{2}\} \;\amalg\;
  \bigl(\mathfrak{S}_{3} \wedge_{\{1,\sigma_{1}\}} E^{*}_{1}\bigr),
\]
c'est-à-dire l'ensemble induit de $E^{*}_{1}$, muni de l'action triviale de
$\{1,\sigma_{1}\}$, à $\mathfrak{S}_{3}$.\footnote{Le symbole $\wedge$ est celui
de l'auteur pour l'induction ; il l'emploie aussi, quelques pages plus loin, pour
le produit amalgamé de groupes. Ce sont deux opérations différentes sous un même
signe, et nous ne conservons le signe que là où le manuscrit l'écrit.} On pose
alors $E^{*}_{2} = E_{2} \setminus E'_{2}$, sur lequel la restriction
$p^{*}$ de $p$ satisfait encore $p^{*}\sigma_{3} = \sigma\,p^{*}$.

\subsection*{Quand $E_{2}$ recouvre-t-il $E_{1} \times E_{1}$ ?}

La question posée est celle de la surjectivité de
$(p,\ p\rho) \colon E_{2} \to E_{1} \times E_{1}$ — la même condition de
concaténation que le lot précédent rencontrait sous la forme $(**)$. Sur
$E'_{2}$ l'image se calcule exactement :
\[
  (p, p\rho)(e_{2}) = (e_{1},e_{1}), \qquad
  (p, p\rho)\bigl(\rho^{i}d(\xi)\bigr) =
  \begin{cases}
    (\xi,\ e_{1}) & i = 0, \\
    (e_{1},\ \sigma\xi) & i = 1, \\
    (\sigma\xi,\ \xi) & i = 2,
  \end{cases}
\]
pour $\xi \in E^{*}_{1}$. L'image est donc
$\Sigma = \Sigma_{1} \cup \Sigma_{2} \cup \Delta'$, où $\Sigma_{1} = E_{1}
\times \{e_{1}\}$, $\Sigma_{2} = \{e_{1}\} \times E_{1}$ et $\Delta'$ est
l'image de $\delta' = (\sigma, \mathrm{id})$ ; les trois se coupent deux à deux,
et toutes trois ensemble, en le seul point $(e_{1},e_{1})$.

D'où la réponse, complète : $\Sigma = E_{1} \times E_{1}$ si et seulement si
$E_{1}$ a au plus deux éléments. En effet le complémentaire de
$\Sigma_{1} \cup \Sigma_{2}$ est $E^{*}_{1} \times E^{*}_{1}$, et il faut que
$\xi \mapsto (\sigma\xi,\xi)$ y soit surjectif, c'est-à-dire que
$\eta = \sigma\xi$ pour tous $\xi,\eta \in E^{*}_{1}$ ; si $E^{*}_{1}$ avait au
moins deux éléments, deux orbites distinctes sous $\sigma$, ou une orbite à deux
éléments, fourniraient un contre-exemple.

\subsection*{Le carré cartésien}

Il reste alors à examiner le cas $E_{1} = \{e_{1},\xi\}$, où $E_{2} = E'_{2}$ a
quatre éléments. La question est de savoir si le carré
\begin{tikzcd}[column sep=large, row sep=small]
E_{1} & E_{2} \arrow[l, "p"'] \\
E_{0} \arrow[u] & E_{1} \arrow[l] \arrow[u, "\delta_{3}"']
\end{tikzcd}
est cartésien, c'est-à-dire si $p^{-1}(e_{1}) = \delta_{3}(E_{1})$. L'inclusion
$\supset$ est acquise ; l'autre se vérifie sur la classification ci-dessus. Un
élément de $E'_{2}$ distinct de $e_{2}$ s'écrit $\rho^{i}d(\eta)$ avec
$\eta \neq e_{1}$, et $p$ y vaut $\eta$ si $i = 0$, $\sigma\eta$ si $i = 2$ —
tous deux distincts de $e_{1}$. Donc $p(\zeta) = e_{1}$ force $i = 1$, et
$\rho\,d(E_{1}) = \delta_{3}(E_{1})$ : le carré est bien cartésien.

L'identification $\delta_{3} = \rho\,d\sigma$ qui sert ici est écrite sur la
page avec l'exposant de $\rho$ raturé ; le calcul la donne sans ambiguïté,
puisque $d\sigma(x,y) = (y,x,x)$ et $\rho(y,x,x) = (x,x,y) =
\delta_{3}(x,y)$.\footnote{C'est le seul endroit du dossier où un exposant
illisible se laisse restituer par le calcul plutôt que deviné. La transcription
laisse le trou ; nous le comblons, et disons ici par quoi.}

\pagerange{26}{29}
\section*{1. La situation type : un fermé dans un espace}

Le texte annoncé par le sommaire de la page 1 commence ici. La donnée est
\[
  Y \hookrightarrow X, \qquad Y \text{ fermé dans } X,
  \qquad X^{*} = X \setminus Y,
\]
des espaces topologiques, avec l'avertissement qu'il faudra plus tard les
remplacer par des topos.

\subsection*{Trivialité normale locale}

L'hypothèse posée est que la « structure normale » de $X$ le long de $Y$ soit
localement triviale : pour tout $y \in Y$, il existe un voisinage ouvert $X'$ de
$y$ dans $X$, un espace $Z$ pointé par $z$, et un homéomorphisme
$X' \simeq Y' \times Z$ (où $Y' = Y \cap X'$) échangeant l'inclusion
$Y' \hookrightarrow X'$ et $Y' \times \{z\} \hookrightarrow Y' \times Z$ :

\begin{tikzcd}[column sep=large]
Y' \arrow[r, hook] \arrow[d] & X' \arrow[d] \\
Y'\times\{z\} \arrow[r, hook] & Y'\times Z
\end{tikzcd}

L'espace $Z$ est le modèle transverse, et $z$ y est le point singulier : si $Y$
et $X^{*}$ sont des variétés, $Z \setminus \{z\}$ en est une aussi, et $X$ est
« singulier le long de $Y$ ».

\subsection*{Le tube}

On introduit une factorisation canonique de l'inclusion,
\[
  Y \hookrightarrow \mathcal{V}_{Y,X} \hookrightarrow X,
\]
où $\mathcal{V}_{Y,X}$ est le \emph{voisinage tubulaire} — un objet
infinitésimal, dont le manuscrit dit qu'il faudra donner une définition
techniquement souple, et dont il propose, pour $X$ paracompact, le germe de $X$
autour de $Y$ : la limite projective, comme topos, des voisinages ouverts de $Y$
dans $X$.\footnote{Le manuscrit dit que ce germe « est un topos », ce qui est
exact au sens où la limite projective d'un système de topos ouverts en est un ;
c'est aussi ce qui rend la notion utilisable dans le contexte étale, où il n'y a
pas de voisinage ouvert assez petit. La construction reste à cet endroit
programmatique, et le texte y insiste.}

Les deux inclusions ont des comportements opposés, qui sont exactement ce dont
on se servira : $\mathcal{V}_{Y,X} \hookrightarrow X$ a les propriétés d'un
ouvert, et $Y \hookrightarrow \mathcal{V}_{Y,X}$ celles d'une immersion fermée.
Sous les hypothèses de trivialité normale locale — et des conditions de
modération ensemblistes, que le manuscrit invoque sans les
détailler\footnote{Elles sont indispensables : sans elles, le système des
voisinages ouverts peut n'être cofinal dans aucun système raisonnable et la
rétraction n'existe pas. Le texte le sait et l'écrit en marge : l'intuition est
que $\mathcal{V}_{Y,X}$ et $Y$ sont homotopiquement équivalents indépendamment de
la trivialité locale.} — on a une équivalence d'homotopie
\[
  Y \xrightarrow{\ \sim\ } \mathcal{V}_{Y,X}.
\]

Le tube épointé est ce qu'il reste quand on ôte $Y$ :
\[
  \mathcal{V}^{*}_{Y,X} = \mathcal{V}_{Y,X} \cap X^{*}
  = \mathcal{V}_{Y,X} \setminus Y,
\]
et le carré des quatre espaces s'écrit, la flèche du haut étant ce que le
manuscrit appelle joliment le « bouchage de trous le long de $Y$ » :

\begin{tikzcd}[column sep=large]
\mathcal{V}^{*}_{Y,X} \arrow[r, hook] \arrow[d, hook] & \mathcal{V}_{Y,X} \arrow[d, hook] \\
X^{*} \arrow[r, hook] & X
\end{tikzcd}

\subsection*{$X$ comme somme amalgamée}

L'assertion, posée indépendamment de toute hypothèse de trivialité ou de
lissité, est que $X$ s'obtient par recollement de $\mathcal{V}_{Y,X}$ et de
$X^{*}$ le long de $\mathcal{V}^{*}_{Y,X}$ : c'est la somme amalgamée du
diagramme
\[
  X^{*} \longleftarrow \mathcal{V}^{*}_{Y,X} \hookrightarrow \mathcal{V}_{Y,X}.
\]
Au niveau des espaces, l'énoncé est vrai et élémentaire dès que le tube est
remplacé par un vrai voisinage ouvert : si $U$ est un ouvert contenant $Y$,
alors $\{U, X^{*}\}$ est un recouvrement ouvert de $X$, $U \cap X^{*} = U
\setminus Y$, et tout espace est la somme amalgamée de deux ouverts le long de
leur intersection.\footnote{C'est l'hypothèse que le manuscrit ne formule pas et
qui donne à l'énoncé son contenu : la somme amalgamée est prise dans les
espaces, pour le recouvrement ouvert $\{U,X^{*}\}$, et le passage au germe
$\mathcal{V}_{Y,X}$ est un passage à la limite sur $U$. C'est aussi la raison
pour laquelle il tient à ce que $\mathcal{V}_{Y,X} \hookrightarrow X$ « ait les
propriétés d'un ouvert ».} Ce que le manuscrit veut, et qui demande beaucoup
plus, est la version homotopique : que le type d'homotopie de $X$ soit la somme
amalgamée \emph{homotopique} des types d'homotopie des trois autres.

\pagerange{29}{33}
\section*{Van Kampen, et la fibration du tube épointé}

\subsection*{Le carré cocartésien de groupoïdes}

Le premier étage du programme est celui des groupoïdes fondamentaux. Le carré

\begin{tikzcd}[column sep=large]
\Pi_{1}\mathcal{V}^{*}_{Y,X} \arrow[r, "p"] \arrow[d, "i"'] & \Pi_{1}\mathcal{V}_{Y,X} \arrow[d, dashed] \\
\Pi_{1}X^{*} \arrow[r, dashed] & \Pi_{1}X
\end{tikzcd}

doit être cocartésien.\footnote{Les deux flèches qui le ferment sont dessinées
en pointillé sur la page, et les deux termes de droite mis entre parenthèses :
c'est ce que l'énoncé doit fournir, non ce qu'on tient déjà. Nous conservons le
pointillé, qui porte cette distinction.} C'est le théorème de van Kampen sous sa
forme groupoïdale, celle qui n'exige ni connexité ni choix de point base : pour
un recouvrement ouvert par deux pièces, le groupoïde fondamental du tout est la
somme amalgamée des groupoïdes fondamentaux des pièces au-dessus de celui de
l'intersection.\footnote{La forme groupoïdale est due à R.~Brown (1967) ; elle
est antérieure à ces pages, et c'est précisément elle qui rend l'énoncé
utilisable ici, où $\mathcal{V}^{*}_{Y,X}$ n'a aucune raison d'être connexe. Le
manuscrit passe ensuite à la forme pointée, qui suppose $Y$, $X$, $X^{*}$
connexes et un point base choisi — un revêtement universel de
$\mathcal{V}^{*}_{Y,X}$, dit-il, ce qui est en effet la bonne façon de pointer
un groupoïde.}

Sous ces hypothèses de connexité, l'énoncé prend la forme pointée
\[
  \Pi_{1}(X) \;\simeq\;
  \Pi_{1}(X^{*}) \ast_{\Pi_{1}\mathcal{V}^{*}_{Y,X}} \Pi_{1}(\mathcal{V}_{Y,X}),
  \qquad \Pi_{1}(\mathcal{V}_{Y,X}) \simeq \Pi_{1}(Y),
\]
un produit amalgamé de groupes. Tout revient donc à expliciter
$\Pi_{1}(\mathcal{V}^{*}_{Y,X})$ et ses deux images, c'est-à-dire le diagramme

\begin{tikzcd}[column sep=large, row sep=small]
\Pi_{1}(\mathcal{V}^{*}_{Y,X}) \arrow[r] \arrow[d] & \Pi_{1}(Y) \\
\Pi_{1}(X^{*}) &
\end{tikzcd}

\subsection*{Ce que sert la trivialité normale locale}

C'est ici, et seulement ici, que l'hypothèse de trivialité normale locale
travaille : elle doit fournir que l'inclusion
$\mathcal{V}^{*}_{Y,X} \hookrightarrow \mathcal{V}_{Y,X}$ se comporte comme une
\emph{fibration homotopique}, de base $Y \simeq \mathcal{V}_{Y,X}$ et de fibre
le modèle transverse épointé $\mathcal{V}^{*}_{z,Z}$ — le \emph{lien} de la
strate, en langue d'aujourd'hui. Le manuscrit propose de faire de cette
propriété la \emph{définition} de la trivialité normale locale dans les
contextes où la formulation naïve n'a pas de sens, celui des schémas pour la
topologie étale au premier chef ; l'hypothèse de lissité devient alors
techniquement inutile et n'était qu'heuristique.\footnote{La fibre est écrite
$\mathcal{V}^{*}_{z,Z}$ dans la phrase et $\mathcal{V}^{*}_{z,X}$ dans la
formule deux lignes plus bas. C'est $Z$, le modèle transverse, qui est le bon :
$z$ n'est pas un point de $X$. Dans le contexte schématique, ajoute-t-il, il
faudra travailler avec des types d'homotopie profinis et les localiser hors des
caractéristiques résiduelles, comme chez Artin--Mazur.}

D'où la suite exacte d'homotopie de la fibration,
\[
  \cdots \to \Pi_{i}(\mathcal{V}^{*}_{z,Z}) \to
  \Pi_{i}(\mathcal{V}^{*}_{Y,X}) \to \Pi_{i}(Y) \to
  \Pi_{i-1}(\mathcal{V}^{*}_{z,Z}) \to \cdots
\]
qui n'intéresse qu'en basses dimensions. Si la fibre $\mathcal{V}^{*}_{z,Z}$ est
connexe, la flèche $\Pi_{1}(\mathcal{V}^{*}_{Y,X}) \to \Pi_{1}(Y)$ est
surjective et $\Pi_{1}(\mathcal{V}^{*}_{Y,X})$ est une extension de $\Pi_{1}(Y)$
par le quotient de $\Pi_{1}(\mathcal{V}^{*}_{z,Z})$ par l'image de
$\Pi_{2}(Y)$ — donc par $\Pi_{1}(\mathcal{V}^{*}_{z,Z})$ lui-même si
$\Pi_{2}(Y)$ est nul.\footnote{L'hypothèse est écrite sur la page sous la forme
« si $z$ est déconnectant dans $Z$ localement, i.e. si la fibre
$\mathcal{V}^{*}_{z,Z}$ est connexe » ; le mot est signalé comme lecture
incertaine par la transcription, et il dit en français le contraire de la glose
qui le suit — un point qui déconnecte laisse un lien \emph{non} connexe, comme
un hyperplan dans $\mathbb{R}^{n}$. C'est la glose qui porte l'argument, et
c'est elle que nous retenons : la fibre est supposée connexe.}

Lorsque de plus $\Pi_{2}(Y) \to \Pi_{1}(\mathcal{V}^{*}_{z,Z})$ est nulle, on
obtient une extension
\[
  1 \longrightarrow \Pi_{1}(\mathcal{V}^{*}_{z,Z}) \longrightarrow
  \Pi_{1}(\mathcal{V}^{*}_{Y,X}) \longrightarrow \Pi_{1}(Y)
  \longrightarrow 1,
\]
accompagnée de la flèche $\Pi_{1}(\mathcal{V}^{*}_{Y,X}) \to \Pi_{1}(X^{*})$.
Grothendieck baptise le sous-groupe de gauche \emph{inertie} $I$ et le groupe du
milieu \emph{décomposition} $D$ : les noms sont ceux de la théorie des nombres,
et l'analogie est exacte — l'inertie est la monodromie locale autour de la
strate, la décomposition tout ce que le tube épointé
voit.\footnote{Le manuscrit porte les deux noms sous la suite sans dire
explicitement auquel des termes le second se rapporte ; la transcription le
signale. Le rapprochement avec le groupe de décomposition d'une place fixe la
lecture : $D$ est le groupe fondamental du tube épointé, $I$ le noyau de sa
surjection sur celui de la strate.}

Un cas est distingué, et c'est celui dont tout le dossier se sert : celui où les
pièces $Y$, $X^{*}$, $\mathcal{V}^{*}_{z,Z}$ sont des $K(\pi,1)$, c'est-à-dire
ont tous leurs $\Pi_{i}$ nuls pour $i \geq 2$. La suite exacte montre alors
qu'il en est de même de $\mathcal{V}^{*}_{Y,X}$, et bien sûr de
$\mathcal{V}_{Y,X} \simeq Y$.

\pagerange{34}{37}
\section*{Les faisceaux constructibles : la catégorie $\mathfrak{F}$}

Indépendamment de toute hypothèse de nullité ou de trivialité locale, une chose
se reconstruit à partir du seul carré de groupoïdes : la catégorie
$\mathfrak{F}$ des faisceaux $F$ sur $X$ tels que $F|X^{*}$ et $F|Y$ soient
localement constants. Ce sont les faisceaux qu'on appelle aujourd'hui
\emph{constructibles} pour la stratification à deux strates.

Cette catégorie est équivalente à celle des triplets
\[
  \bigl(E_{X^{*}},\ E_{Y,X}\,;\ \varphi\bigr),
\]
où $E_{X^{*}}$ est un système local sur $\Pi_{1}X^{*}$, $E_{Y,X}$ un système
local sur $\Pi_{1}\mathcal{V}_{Y,X} \simeq \Pi_{1}Y$ — c'est-à-dire un
revêtement étale de $Y$ —, et où $\varphi$ est un morphisme de systèmes locaux
sur $\Pi_{1}\mathcal{V}^{*}_{Y,X}$ :
\[
  \varphi \colon p^{*}(E_{Y,X}) \longrightarrow i^{*}(E_{X^{*}}).
\]
Lorsque les trois groupoïdes sont connexes, cela se dit sans groupoïdes :
$E_{X^{*}}$ est un ensemble à opérateurs $\Pi_{1}(X^{*})$, $E_{Y,X}$ un
ensemble à opérateurs $\Pi_{1}(Y)$, et $\varphi$ une application entre eux
compatible aux opérations restreintes à $\Pi_{1}(\mathcal{V}^{*}_{Y,X})$.

Le manuscrit en tire que $\mathfrak{F}$ est un topos de préfaisceaux,
\[
  \mathfrak{F} \;\simeq\; \widehat{\mathcal{C}},
\]
sur une catégorie $\mathcal{C}$ fabriquée en collant les trois groupoïdes : ses
flèches sont celles des trois groupoïdes, plus, d'un objet $a$ de
$A = \Pi_{1}\mathcal{V}^{*}_{Y,X}$ vers un objet $b$ de
$B = \Pi_{1}\mathcal{V}_{Y,X}$ ou $c$ de $C = \Pi_{1}X^{*}$, les flèches
$p^{*}(a) \to b$ de $B$, respectivement $i^{*}(a) \to c$ de $C$. C'est le
\emph{collage} des deux foncteurs — le cographe du profoncteur de spécialisation
— et c'est, en langue d'aujourd'hui, la troncation en dimension~1 de la
\emph{catégorie des chemins sortants} de l'espace
stratifié.\footnote{La page 37 est recouverte de deux tiers de ratures et la
phrase qui définit $\mathcal{C}$ y est reprise deux fois, avec un mot illisible
au point critique. Nous ne prétendons pas restituer la définition mot pour mot ;
ce qui est sûr, et ce dont les deux calculs qui suivent se servent, est que
$\mathcal{C}$ est obtenue en collant $B$ et $C$ le long de $A$, et que
$\mathfrak{F}$ en est les préfaisceaux.}

\subsection*{Le plan complexe}

Le cas typique pour la géométrie algébrique est celui d'une surface topologique
$X$ et d'un point $Y = \{a\}$ : le tube est le disque infinitésimal $D_{a}$, le
tube épointé le disque épointé $D^{*}_{a}$, et le diagramme des $\Pi_{1}$ se
réduit à

\begin{tikzcd}[column sep=large, row sep=small]
\mathbb{Z}_{a,X} \arrow[r] \arrow[d] & \{1\} \\
\Pi_{1}(X^{*}) &
\end{tikzcd}

où $\mathbb{Z}_{a,X} \simeq \mathbb{Z}$ est le module d'orientation de $X$ en
$a$, dont les deux générateurs correspondent aux deux orientations, et où la
flèche verticale est celle qui définit la structure locale en $a$ du groupe
fondamental de $X^{*}$. Le produit d'un nombre arbitraire de copies de ce modèle
donne les modèles-types des stratifications associées aux diviseurs à croisements
normaux.

Pour $X = \mathbb{C}$ et $a = 0$, la flèche $A \to C$ est un isomorphisme
$\Pi_{1}(D^{*}_{0}) \simeq \Pi_{1}(\mathbb{C}^{*})$ : la catégorie
$\mathcal{C}$ se réduit à deux objets, $\mathrm{Aut}(a) = \mathbb{Z}$ et une
seule flèche $a \to b$, de sorte que $b$ y est objet final. Une catégorie à
objet final a un nerf contractile ; donc $\widehat{\mathcal{C}}$, et avec lui
$\mathfrak{F}$, est de type d'homotopie trivial — ce qui est bien le type
d'homotopie de $X = \mathbb{C}$.

\pagerange{37}{40}
\section*{La sphère, ou pourquoi le niveau $\Pi_{1}$ ne suffit pas}

Prenons maintenant le compactifié :
\[
  X = \mathbb{P}^{1}_{\mathbb{C}} \simeq S^{2}, \qquad
  Y = \{\infty\}, \qquad X^{*} = \mathbb{C}.
\]
Le carré de recollement est

\begin{tikzcd}[column sep=large]
D^{*}_{\infty} \arrow[r, hook] \arrow[d] & D_{\infty} \arrow[d, dashed] \\
\mathbb{C} \arrow[r, dashed] & \mathbb{P}^{1}_{\mathbb{C}} \simeq S^{2}
\end{tikzcd}

et sa somme amalgamée est bien $S^{2}$, dont le type d'homotopie n'a rien de
trivial. Le diagramme correspondant sur les $\Pi_{1}$, lui, est

\begin{tikzcd}[column sep=large, row sep=small]
\mathbb{Z} \arrow[r] \arrow[d] & 1 \\
1 &
\end{tikzcd}

et la catégorie $\mathcal{C}$ qui en sort est la catégorie
${\cdot} \to {\cdot}$, à deux objets et une flèche. Elle a un objet final, donc
$\mathfrak{F} \simeq \widehat{\mathcal{C}}$ est \emph{encore homotopiquement
trivial}, alors que $X$ ne l'est pas — et les pièces de construction
$\{\infty\}$, $\mathbb{C}$, $D^{*}_{\infty}$ sont pourtant toutes des
$K(\pi,1)$, les deux premières étant même contractiles. Le morphisme canonique
\[
  X \longrightarrow \mathfrak{F}
\]
n'est donc pas une équivalence d'homotopie.

\subsection*{Ce que la troncation a perdu}

Le diagnostic se dit en une phrase aujourd'hui, et il vaut la peine de le dire,
parce qu'il montre que rien n'était faux dans la construction — seulement trop
grossier. La catégorie des chemins sortants de $\mathbb{P}^{1}$ stratifié par
$\{\infty\}$ a deux objets, et l'espace des chemins allant de $\infty$ vers le
stratum ouvert est le lien, $D^{*}_{\infty} \simeq S^{1}$. L'espace classifiant
de cette catégorie \emph{à homotopie près} est la somme amalgamée homotopique
$\ast \leftarrow S^{1} \to \ast$, c'est-à-dire $S^{2}$ : rien n'est perdu.
Tronquer en dimension~1, c'est remplacer cet espace de chemins $S^{1}$ par son
$\pi_{0}$, un point — d'où la flèche unique $a \to b$ de $\mathcal{C}$, d'où
l'objet final, d'où la contractilité. C'est exactement ce cran-là que le calcul
de la page 39 met en évidence.\footnote{La formulation moderne — faisceaux
constructibles contre représentations de la catégorie des chemins sortants,
et espace classifiant de celle-ci contre type d'homotopie de l'espace — est
postérieure de plus de vingt ans à ces pages (MacPherson, Treumann, Lurie,
Ayala--Francis--Tanaka). Ce que ces trois pages contiennent est l'observation qui
la rend nécessaire : la troncation en dimension~1 ne suffit pas, et la
raison n'est pas que les pièces soient compliquées, puisqu'elles sont ici les
plus simples possibles.}

Le manuscrit tire lui-même la bonne conclusion, et c'est sur cette phrase que le
lot s'achève :

\begin{quote}
J'avais craint un moment, par cet exemple, que le type d'homotopie de $X$ ne
pouvait pas se récupérer à partir de ceux de $X^{*}$, $Y$,
$\mathcal{V}^{*}_{Y,X}$ et du diagramme de recollement. Mais je me suis
convaincu ensuite que cette crainte n'était pas fondée : le type $\mathfrak{F}$
n'est pas suffisant, lui, pour exprimer le type d'homotopie.
\end{quote}

Ce qui échoue est un invariant, non le programme. Reste à trouver l'invariant
qui marche, et l'essai commence aussitôt, sur le cas de $S^{2}$ : récupérer
$H^{2}(X,A)$ en interprétant cette cohomologie comme classifiant les
\emph{gerbes} sur $X$ liées par $A$. La phrase s'interrompt au bas de la page
40 ; elle se termine à la page 41, dans le lot suivant, et le calcul
réussit.\footnote{Le texte s'arrête au milieu d'une phrase, en bas de feuillet.
Ce que la phrase allait dire n'est pas deviné ici : il est lu à la page 41, qui
la complète.}

\end{document}
